\documentclass[10pt,a4paper]{amsart}
\usepackage[margin=19mm]{geometry}
\usepackage{amsmath,amssymb,mathtools}
\usepackage[scr=rsfs,cal=euler]{mathalfa}
\usepackage{xfrac}
\usepackage[unicode,hidelinks,bookmarks=false]{hyperref}
\newcommand{\ie}{i.e.\ }
\newcommand{\cf}{cf.\ }
\newcommand{\resp}{respectively\ }
\newcommand{\Subsection}{\subsection}
\providecommand{\nobreakdash}{\nobreak\hbox{-}}
\setlength{\emergencystretch}{2em}
\multlinegap0pt
\usepackage{bm}
\usepackage{bbm}
\usepackage{dsfont}
\usepackage{xspace}
\usepackage{cancel}
\usepackage{mathtools}
\usepackage{amsthm}
\makeatletter
\@ifundefined{theorem}{\newtheorem{theorem}{Theorem}}{}
\makeatother
\title[Dynamic Decision Modeling for Viable Short and Long Term Pro]{Dynamic Decision Modeling for Viable Short and Long Term Production Policies: An HJB Approach}
\author{Achraf Bouhmady, Mustapha Serhani, Nadia Raissi}
\date{Source: arXiv:2509.12205v1 (CC BY 4.0)}
\begin{document}
\maketitle

\noindent\emph{Source and licence.} Dynamic Decision Modeling for Viable Short and Long Term Production Policies: An HJB Approach. Version arXiv:2509.12205v1 (CC BY 4.0), distributed under the Creative Commons Attribution 4.0 International licence, as recorded in the arXiv licence metadata for this exact version and shown on the abstract page https://arxiv.org/abs/2509.12205v1. Full rights notice: licenses/arxiv-2509.12205-CC-BY-4.0.txt.\par\smallskip

\noindent\textit{Editorial scope.} Counted unit: the Theorem whose source label is \texttt{thm1} in the article (source0.tex lines 375--381, source label \texttt{thm1}), delivered together with the article's own complete original proof of it (source0.tex lines 384--439, one \texttt{proof} environment of 56 lines). No mathematics is written, repaired or re-derived: statement and proof are the article's own text line for line. Required context retained uncounted: utility (lines 342--344). Every definition, display and label of those lines is retained unchanged. Omitted from this package and not redistributed: 1--341, 345--374, 382--383, 440--1364. Nothing inside a retained excerpt is omitted or altered: every retained body is delivered exactly as the article's lines are written, so no omission receipt of any kind accompanies this package. Citations: the retained excerpts carry no citation command at all, so no bibliography entry of the article is transcribed and no reference is redistributed. Reference adaptation: none was needed and none was made; every reference inside the delivered unit resolves inside it, so no unresolved reference or citation is produced. Printed slips kept literal: no printed word, symbol, hypothesis or number of the counted statement or of its proof was altered. Layout and class: the article's own class is \texttt{interact} with packages including epstopdf, subfig; the wrapper is the lane's pinned \texttt{amsart} editorial wrapper at 10pt on A4, which restates the theorem-like environments the retained text needs and adds the article's own macro definitions that the retained text needs (none), with extra wrapper packages (bm, bbm, dsfont, xspace, cancel, mathtools). The delivered numbering is the unit's own. Assets: no retained span contains a figure environment, a graphics-inclusion command, a PDF-page inclusion, a table or a TikZ picture; no image file is copied into this package, the package declares no asset dependency and no image file is redistributed with it. Licence: version arXiv:2509.12205v1 of this article is distributed under the Creative Commons Attribution 4.0 International licence, as recorded in the arXiv licence metadata for this exact version and shown on the abstract page https://arxiv.org/abs/2509.12205v1. A full rights notice is delivered at licenses/arxiv-2509.12205-CC-BY-4.0.txt. Characters: every retained excerpt is pure ASCII, so the delivered engine, XeLaTeX with its default Latin Modern text font, reports no missing character.
\begin{equation}\label{utility}
v(t, \mathbf{y}) = \inf_{\omega \in \hat{\mathcal{U}}}\left\{ \max \left(  v_0(\mathbf{X}_{\mathbf{y}}^\omega(T)), \sup_{\theta \in (t, T)} g\left(\mathbf{X}_{\mathbf{y}}^\omega(\theta)\right) \right)  \right\},
\end{equation}

\begin{theorem}\label{thm1}
%Let \( v_0 \) and \( g \) be Lipschitz continuous functions defined by (\ref{v01}) and (\ref{g1}), respectively, since $K$ and $C$ are closed.
 Let \( v \) be the value function defined in (\ref{utility}). Then, for every \( t \geq 0 \), the capture basin is given by
\[ 
\operatorname{Cap}_{c}(t) = \{\mathbf{y} \in \Omega \mid v(t, \mathbf{y}) \leq 0\}.
\]
\end{theorem}

\begin{proof}
Assume \( v(t, \mathbf{y}) \leq 0 \). Given that the dynamic \( f(\cdot,\omega) \) is Lipschitz continuous for fixed $\omega\in \mathbb{U}$, and \( f(x, \mathbb{U}) \) is convex for each $x$, an admissible trajectory \( \mathbf{X}_{\mathbf{y}}^\omega \) exists such that
\[
\max \left( v_0(\mathbf{X}_{\mathbf{y}}^\omega(T)), \sup_{\theta \in (t, T)} g\left(\mathbf{X}_{\mathbf{y}}^\omega(\theta)\right) \right) \leq 0.
\]

Therefore, for all \( \theta \in (t, T) \),
\[
g\left(\mathbf{X}_{\mathbf{y}}^\omega(\theta)\right) \leq 0,
\]
This implies that \( \mathbf{X}_{\mathbf{y}}^\omega(\theta) \in \hat{K} \) for all \( \theta \in (t, T) \).

Additionally,
\[
v_0(\mathbf{X}_{\mathbf{y}}^\omega(T)) \leq 0,
\]
This implies that \( \mathbf{X}_{\mathbf{y}}^\omega(T) \in C \). Thus, this defines \( \operatorname{Cap}_{C}(t) \).\\
To prove the other inclusion, assume \( \mathbf{y} \in \operatorname{Cap}_{C}(t) \). By the definition of the capture basin \( \operatorname{Cap}_{C}(t) \), there exists an admissible control \( \omega \in \hat{\mathcal{U}} \) and an associated trajectory \( \mathbf{X}_{\mathbf{y}}^\omega(s) \) for \( s \in [t, T] \) such that

For all \( s \in (t, T) \),
\[
\mathbf{X}_{\mathbf{y}}^\omega(s) \in \hat{K},
\]
which implies that
\[
g\left(\mathbf{X}_{\mathbf{y}}^\omega(s)\right) \leq 0.
\]

At the terminal time \( T \),
\[
\mathbf{X}_{\mathbf{y}}^\omega(T) \in C,
\]
which implies that
\[
v_0\left(\mathbf{X}_{\mathbf{y}}^\omega(T)\right) \leq 0.
\]

Now, consider the value function \( v(t, \mathbf{y}) \), defined by
\[
v(t, \mathbf{y}) = \inf_{\omega \in \hat{\mathcal{U}}} \max \left( v_0\left(\mathbf{X}_{\mathbf{y}}^\omega(T)\right), \sup_{s \in (t, T)} g\left(\mathbf{X}_{\mathbf{y}}^\omega(s)\right) \right).
\]
Since we have an admissible control \( \omega \) such that
\[
\max \left( v_0\left(\mathbf{X}_{\mathbf{y}}^\omega(T)\right), \sup_{s \in (t, T)} g\left(\mathbf{X}_{\mathbf{y}}^\omega(s)\right) \right) \leq 0,
\]
it follows that
\[
v(t, \mathbf{y}) \leq 0.
\]

This establishes the other inclusion, so we conclude that
\[
\operatorname{Cap}_{C}(t) = \left\{ \mathbf{y} \in \Omega \mid v(t, \mathbf{y}) \leq 0 \right\},
\]
This completes the proof of the theorem.
\end{proof}

\end{document}
